\documentclass[11pt,a4paper]{article}

% Pakete
\usepackage[utf8]{inputenc}
\usepackage[T1]{fontenc}
\usepackage[ngerman]{babel}
\usepackage[left=18mm,right=18mm,top=18mm,bottom=18mm]{geometry}
\usepackage{helvet}
\renewcommand{\familydefault}{\sfdefault}
\usepackage{amsmath,amssymb}
\usepackage{graphicx}
\usepackage{tikz}
\usetikzlibrary{arrows.meta,positioning,patterns}
\usepackage{tcolorbox}
\tcbuselibrary{skins,breakable}
\usepackage{enumitem}
\usepackage{tabularx}
\usepackage{colortbl}
\usepackage{qrcode}
\usepackage{xcolor}
\usepackage[hidelinks]{hyperref}

% Fachfarbe Physik
\definecolor{physik}{HTML}{2c5aa0}
\definecolor{physiklight}{HTML}{e8f0fa}

% Box-Stile
\tcbset{
    kasten/.style={
        colback=physiklight,
        colframe=physik,
        fonttitle=\bfseries,
        title=#1,
        boxrule=1pt,
        arc=2pt,
        left=5pt,
        right=5pt,
        top=5pt,
        bottom=5pt
    },
    formelbox/.style={
        colback=white,
        colframe=physik,
        boxrule=1pt,
        arc=2pt,
        left=6pt,
        right=6pt,
        top=6pt,
        bottom=6pt
    }
}

% Lücke
\newcommand{\luecke}[1]{\underline{\hspace{#1}}}

% Kopfzeile
\usepackage{fancyhdr}
\pagestyle{fancy}
\fancyhf{}
\lhead{\small Physik 12 GK}
\chead{\small Aufenthaltswahrscheinlichkeit}
\rhead{\small Name: \underline{\hspace{4cm}}}
\renewcommand{\headrulewidth}{0.4pt}

% Keine Einrückung
\setlength{\parindent}{0pt}
\setlength{\parskip}{6pt}

\begin{document}

% Titel
\begin{minipage}[t]{0.65\textwidth}
\vspace{0pt}
{\Large\textbf{Aufenthaltswahrscheinlichkeit}}\\[0.3em]
{\normalsize Fr 23.01.2026}
\end{minipage}
\hfill
\begin{minipage}[t]{0.32\textwidth}
\vspace{0pt}
\raggedleft
\begin{tabular}{@{}c@{\hspace{8pt}}c@{}}
\href{https://devbox42.github.io/sciencesim/physik/kl12-GK/01-photoeffekt/05-DS-uebung/LP-05-wahrscheinlichkeit.html}{\qrcode[height=1.1cm]{https://devbox42.github.io/sciencesim/physik/kl12-GK/01-photoeffekt/05-DS-uebung/LP-05-wahrscheinlichkeit.html}} &
\href{https://devbox42.github.io/sciencesim/physik/kl12-GK/01-photoeffekt/05-DS-uebung/SIM-05-elektronenbeugung.html}{\qrcode[height=1.1cm]{https://devbox42.github.io/sciencesim/physik/kl12-GK/01-photoeffekt/05-DS-uebung/SIM-05-elektronenbeugung.html}} \\
{\scriptsize Lernpfad} & {\scriptsize e-Beugung}
\end{tabular}
\end{minipage}

\vspace{0.3em}
\hrule height 1pt
\vspace{0.8em}

% ============================================================
% KASTEN 1: Stochastische Vorhersagbarkeit
% ============================================================

\begin{tcolorbox}[kasten={Kasten 1: Stochastische Vorhersagbarkeit}]

Ein einzelnes Elektron trifft am Doppelspalt an einem \luecke{2.5cm} Ort auf.

Dieser Ort ist nicht \luecke{2.5cm}.

\vspace{0.5em}

Schickt man jedoch \luecke{2cm} Elektronen durch den Doppelspalt,

entsteht ein \luecke{3cm}. Diese Verteilung ist vorhersagbar.

\vspace{0.8em}

\begin{center}
\fbox{\parbox{0.85\textwidth}{\centering\vspace{0.3em}
\textbf{Merksatz:} \glqq\luecke{2.5cm}, in Summe \luecke{2.5cm}\grqq
\vspace{0.3em}}}
\end{center}

\vspace{0.3em}
{\small\textit{Wortbank: zufällig, vorhersagbar, viele, Interferenzmuster, einzeln zufällig, determiniert}}
\end{tcolorbox}

\vspace{0.6em}

% ============================================================
% KASTEN 2: Wahrscheinlichkeitswelle
% ============================================================

\begin{tcolorbox}[kasten={Kasten 2: Die Wahrscheinlichkeitswelle}]

Die Wellenfunktion $\psi$ beschreibt nicht den genauen Ort eines Elektrons.

Sie beschreibt die \luecke{3cm}, das Elektron an einem bestimmten Ort zu finden.

\vspace{0.5em}

Die Größe $|\psi|^2$ heißt \luecke{5cm}.

\vspace{0.8em}

\textbf{Im Interferenzmuster:}
\begin{itemize}[leftmargin=1.5cm, itemsep=2pt]
\item Helle Streifen: $|\psi|^2$ ist \luecke{2cm} $\rightarrow$ hohe Trefferwahrscheinlichkeit
\item Dunkle Streifen: $|\psi|^2$ ist \luecke{2cm} $\rightarrow$ niedrige Trefferwahrscheinlichkeit
\end{itemize}

\vspace{0.3em}
{\small\textit{Wortbank: groß, klein, Wahrscheinlichkeit, Aufenthaltswahrscheinlichkeitsdichte}}
\end{tcolorbox}

\vspace{1em}

% ============================================================
% AUFGABE 1: Histogramm
% ============================================================

{\large\textbf{Aufgabe 1: Histogramm interpretieren}} \hfill \textbf{(5 BE)}

\vspace{0.5em}

Das Histogramm zeigt die Verteilung von Elektronen am Doppelspalt für unterschiedliche Anzahlen.

\begin{center}
\begin{tikzpicture}[scale=0.9]
    % n = 10
    \begin{scope}[xshift=0cm]
        \draw[thick] (0,0) rectangle (2.5,2.5);
        \node[above] at (1.25,2.5) {\small $n = 10$};
        % Zufällige Balken
        \fill[physik!70] (0.3,0) rectangle (0.5,0.8);
        \fill[physik!70] (0.7,0) rectangle (0.9,0.3);
        \fill[physik!70] (1.1,0) rectangle (1.3,1.2);
        \fill[physik!70] (1.5,0) rectangle (1.7,0.5);
        \fill[physik!70] (1.9,0) rectangle (2.1,0.9);
    \end{scope}

    % n = 100
    \begin{scope}[xshift=3.5cm]
        \draw[thick] (0,0) rectangle (2.5,2.5);
        \node[above] at (1.25,2.5) {\small $n = 100$};
        % Leicht erkennbares Muster
        \fill[physik!70] (0.2,0) rectangle (0.4,0.5);
        \fill[physik!70] (0.5,0) rectangle (0.7,1.2);
        \fill[physik!70] (0.8,0) rectangle (1.0,0.6);
        \fill[physik!70] (1.1,0) rectangle (1.3,1.8);
        \fill[physik!70] (1.4,0) rectangle (1.6,0.7);
        \fill[physik!70] (1.7,0) rectangle (1.9,1.4);
        \fill[physik!70] (2.0,0) rectangle (2.2,0.4);
    \end{scope}

    % n = 1000
    \begin{scope}[xshift=7cm]
        \draw[thick] (0,0) rectangle (2.5,2.5);
        \node[above] at (1.25,2.5) {\small $n = 1000$};
        % Klares Interferenzmuster
        \fill[physik!70] (0.1,0) rectangle (0.25,0.3);
        \fill[physik!70] (0.3,0) rectangle (0.45,1.0);
        \fill[physik!70] (0.5,0) rectangle (0.65,0.2);
        \fill[physik!70] (0.7,0) rectangle (0.85,1.6);
        \fill[physik!70] (0.9,0) rectangle (1.05,0.15);
        \fill[physik!70] (1.1,0) rectangle (1.25,2.2);
        \fill[physik!70] (1.3,0) rectangle (1.45,0.15);
        \fill[physik!70] (1.5,0) rectangle (1.65,1.6);
        \fill[physik!70] (1.7,0) rectangle (1.85,0.2);
        \fill[physik!70] (1.9,0) rectangle (2.05,1.0);
        \fill[physik!70] (2.1,0) rectangle (2.25,0.3);
    \end{scope}
\end{tikzpicture}
\end{center}

\begin{enumerate}[label=\alph*)]
\item Beschreibe, wie sich das Histogramm verändert, wenn die Anzahl der Elektronen steigt. \hfill (2 BE)

\vspace{1.2cm}

\item Erkläre, warum man bei 10 Elektronen noch kein klares Muster erkennt. \hfill (2 BE)

\vspace{1.2cm}

\item Was sagt $|\psi|^2$ über die Orte aus, an denen viele Elektronen auftreffen? \hfill (1 BE)

\vspace{0.8cm}

\end{enumerate}

\newpage

% ============================================================
% AUFGABE 2: Klassisch vs. Quantenmechanisch
% ============================================================

{\large\textbf{Aufgabe 2: Klassisch vs. Quantenmechanisch}} \hfill \textbf{(6 BE)}

\vspace{0.5em}

Vergleiche das Verhalten von Kugeln und Elektronen am Doppelspalt.

\vspace{0.5em}

\begin{tabularx}{\textwidth}{|p{3.5cm}|X|X|}
\hline
\rowcolor{physiklight}
& \textbf{Kugeln (klassisch)} & \textbf{Elektronen (quantenmechanisch)} \\
\hline
\textbf{a) Erwartetes Muster auf dem Schirm}

\vspace{0.3em}
{\small (Skizze)}

\vspace{2.5cm}

& & \\
\hline
\end{tabularx}

\hfill {\small (je 2 BE)}

\vspace{0.8em}

\begin{enumerate}[label=\alph*), start=2]
\item Erkläre den Unterschied zwischen den beiden Mustern. \hfill (2 BE)

\vspace{2cm}

\end{enumerate}

\vspace{0.5em}

% ============================================================
% AUFGABE 3: Elektronenbeugung
% ============================================================

{\large\textbf{Aufgabe 3: Elektronenbeugung}} \hfill \textbf{(6 BE)}

\vspace{0.5em}

Führe das Experiment mit der Elektronenbeugungsröhre durch.

\vspace{0.5em}

\begin{tcolorbox}[formelbox]
\textbf{Theoretische Wellenlänge (de Broglie):} \quad $\lambda_{\text{theo}} = \dfrac{h}{\sqrt{2 \cdot m_e \cdot e \cdot U_B}}$

\vspace{0.3em}
{\small Konstanten: $h = 6{,}626 \cdot 10^{-34}$ J·s, $m_e = 9{,}109 \cdot 10^{-31}$ kg, $e = 1{,}602 \cdot 10^{-19}$ C}
\end{tcolorbox}

\vspace{0.5em}

Miss den Radius des inneren Rings für verschiedene Beschleunigungsspannungen.

\vspace{0.5em}

\begin{center}
\begin{tabular}{|c|c|c|c|c|}
\hline
\rowcolor{physiklight}
$U_B$ (kV) & $\lambda_{\text{theo}}$ (pm) & $r_1$ (cm) & $\lambda_{\text{exp}}$ (pm) & Abweichung (\%) \\
\hline
4,0 & & & & \\[0.8em]
\hline
6,0 & & & & \\[0.8em]
\hline
8,0 & & & & \\[0.8em]
\hline
\end{tabular}
\end{center}

\vspace{0.5em}

\begin{tcolorbox}[formelbox]
\textbf{Auswertungsformel (Kleinwinkelnäherung):} \quad $\lambda_{\text{exp}} \approx \dfrac{d \cdot r_1}{L}$

\vspace{0.3em}
{\small Gerätedaten: $d = 213$ pm (Netzebenenabstand), $L = 13{,}5$ cm (Abstand Folie–Schirm)}
\end{tcolorbox}

\vspace{0.3em}
{\small\textit{Hinweis: $\lambda_{\text{theo}}$ mit de Broglie Formel berechnen, $r_1$ an der Skala ablesen, $\lambda_{\text{exp}}$ mit Auswertungsformel berechnen.}}

\vspace{0.8em}

\begin{enumerate}[label=\alph*)]
\item Wie verändert sich der Ringradius, wenn die Beschleunigungsspannung erhöht wird? \hfill (2 BE)

\vspace{1.2cm}

\item Erkläre in einem Satz, was dieses Experiment über Elektronen zeigt. \hfill (2 BE)

\vspace{1.2cm}

\end{enumerate}

\vspace{1em}

\begin{center}
\rule{10cm}{0.5pt}

\textbf{Gesamt:} \underline{\hspace{1cm}} / 17 BE
\end{center}

\end{document}
